\section{Introduction}

For a positive integer $n$, let $j(n)$ be the least positive integer $m$
such that every interval of $m$ consecutive integers contains an integer
coprime to $n$. Write $\omega(n)$ for the number of distinct prime divisors
of $n$, with $\omega(1)=0$, and put
\[
 h(k)=\sup_{\substack{n\ge1\\\omega(n)\le k}}j(n),\qquad k\ge1.
\]
The starting point of the interval is arbitrary, and replacing $n$ by
the product of its distinct prime divisors does not change $j(n)$.
Thus $h(k)-1$ is the greatest length of an interval that can be covered
by the divisibility classes of at most $k$ primes.

Jacobsthal studied this function in a series of papers beginning in
1960 \cite{Jacobsthal1960}. Erd\H{o}s subsequently recorded Jacobsthal's question
whether $h(k)\ll k^2$ \cite[p.~163, equations (1) and (3)]{Erdos1962}.
In Erd\H{o}s's notation, $C(k)+1$ is the maximum of $j(n)$ over integers
with exactly $k$ distinct prime divisors. This equals our $h(k)$:
adjoining prime divisors cannot decrease $j(n)$.
We prove the following bound, which gives an affirmative answer and a
further iterated-logarithmic saving.

\begin{theorem}\label{thm:main}
There is an absolute constant $C>0$ such that
\[
 h(k)\le C\frac{k^2}{(\log\log(3k))^2}
 \qquad\text{for every integer }k\ge1.
\]
\end{theorem}

All logarithms are natural. In particular the theorem gives a
uniform quadratic bound for arbitrary sets of prime divisors, with
no restriction on their sizes or on the interval's translation.

\subsection*{Historical context and related work}

The development of the upper bound is closely tied to sieve theory.
Erd\H{o}s observed that Brun's method gives $h(k)\ll k^{C_0}$ for some
absolute exponent $C_0$ \cite[p.~163]{Erdos1962}.
Iwaniec's work on the error term in the linear sieve gave an estimate
of order $k^2\log^2 k$ when the prime divisors are the first $k$ primes
\cite{Iwaniec1971}; this primorial case is distinguished explicitly in
\cite[p.~329]{Vaughan1977} and \cite[pp.~225--226]{Iwaniec1978}.
Vaughan then proved the uniform estimate
\[
 j(n)\ll\omega(n)^2\bigl(\log(2\omega(n))\bigr)^4
 \qquad(n>1)
\]
\cite[Theorem, p.~329]{Vaughan1977}.
Iwaniec's shifted-sieve argument removed two logarithmic factors and
established $h(k)\ll k^2\log^2 k$ for arbitrary prime sets
\cite[Theorem and Corollary, p.~226]{Iwaniec1978}.
Theorem~\ref{thm:main} improves this bound and, in particular, reaches
the quadratic scale asked for by Jacobsthal.

Vaughan explained the exponent $2$ through the one-dimensional sieve's
restriction to moduli near the square root of the interval length
\cite[p.~329]{Vaughan1977}. In the usual linear-sieve notation this
appears as the vanishing of the lower function $f$ at sieve parameter
$2$. Our proof also reaches this limiting parameter. It keeps a
boundary contribution that disappears from the leading linear-sieve
term, and then controls the discrepancy between this reference
calculation and the actual prescribed classes. The latter comparison
requires the inverse estimate and the stopped counts developed below;
positivity of the reference calculation alone does not give a lower
bound for the original sieve.

A separate question concerns which prime sets give the largest gap.
Let $p_k$ denote the $k$th prime and $P_k=\prod_{i=1}^k p_i$.
Jacobsthal suggested that $j(P_k)=h(k)$ for every $k$.
Hajdu and Saradha proved equality for $k\le23$, but found
\[
 j(P_{24})=234<236=h(24)
\]
\cite[Theorem~1.2]{HajduSaradha2012}.
Thus the primorial extremality conjecture is distinct from the
quadratic-bound question, and a theorem only about $j(P_k)$ does not
settle the uniform problem. The computational literature reflects
this distinction. Costello and Watts obtained recursive lower bounds
for the minimum number of integers coprime to $P_k$ in an interval,
and used them to compute upper bounds for $j(P_k)$
\cite[Theorem~4.4 and Section~5]{CostelloWatts2015}.
Ziller computed the maximum over arbitrary prime sets through
$43$ prime factors \cite[Section~3]{Ziller2019}.
These finite computations concern explicit values, whereas the
constant in Theorem~\ref{thm:main} is not made explicit.

Primorials nevertheless give the strongest lower bounds relevant to
the order of $h(k)$. They arise from constructions of long covered
intervals, which also produce large gaps between consecutive primes.
Erd\H{o}s's account already places Rankin's lower-bound construction
beside Brun's polynomial upper bound \cite[p.~163, equation (2)]{Erdos1962}.
The later construction of Ford, Green, Konyagin, Maynard and Tao gives
\begin{equation}\label{eq:historical-lower-bound}
 j(P_k)\gg
 k(\log k)^2\frac{\log\log\log k}{\log\log k}
 \qquad(k\longrightarrow\infty)
\end{equation}
\cite[Section~1.1]{FGKMT2018}.
Indeed, their covered-interval bound at prime cutoff $x$ is
\[
 \gg x\log x\,\log\log\log x/\log\log x;
\]
take $x=p_k$ and use
$p_k\sim k\log k$.
Since $j(P_k)\le h(k)$, this supplies a lower bound for the uniform
maximum as well. It is an existence construction and does not assert
the same lower bound for every integer with $k$ prime divisors.
The gap between \eqref{eq:historical-lower-bound} and our upper bound
remains substantial. Vaughan suggested the stronger possible estimate
$j(n)\ll_\eps\omega(n)^{1+\eps}$ for $n>1$ and each fixed $\eps>0$
\cite[p.~329]{Vaughan1977}.

\subsection*{A quantitative covering estimate}

The proof works with one prescribed residue class at each prime below
a size cutoff. This is another form of the same covering problem:
for a finite prime set, the Chinese remainder theorem supplies an
integer $b$ with $b\equiv a_p\pmod p$ at every prime in the set, and
\[
 n\not\equiv a_p\pmod p\ \text{for every such }p
 \quad\Longleftrightarrow\quad
 \gcd\left(n-b,\prod p\right)=1.
\]
This equivalence explains why arbitrary classes and arbitrary
translations occur together; see also
\cite[Lemma~2.1]{HajduSaradha2012}.

For sufficiently large real $z$, set
\begin{equation}\label{eq:parameters}
\begin{gathered}
 L=\log z,\qquad Y=\left\lfloor\frac{z^2}{L^2}\right\rfloor,
 \qquad w=\frac{L}{(\log L)^2},\qquad B=\frac{L}{\log w},\\
 V_0=\prod_{p\le w}\left(1-\frac1p\right).
\end{gathered}
\end{equation}
Products indexed by a letter explicitly called prime always range over
primes. The two cutoffs separate the small primes, which are sieved
inside arithmetic progressions, from the primes in $(w,z]$, which
index the expansion used in the proof. The quantity $B$ is the
logarithmic size of the upper cutoff in base $w$.

\begin{theorem}\label{thm:survivors}
There are absolute constants $c>0$ and $z_0$ with the following property.
For every real $z\ge z_0$ and every choice of a residue class
$a_p\pmod p$ for each prime $p\le z$,
\begin{equation}\label{eq:main-survivors}
 \#\{1\le n\le Y:n\not\equiv a_p\pmod p\text{ for every }p\le z\}
 \ge c\frac{YV_0}{B^2}.
\end{equation}
\end{theorem}

Mertens' theorem identifies the size of this lower bound:
\begin{equation}\label{eq:survivor-scale}
 \frac{YV_0}{B^2}\sim e^{-\gamma}\frac{Y\log w}{L^2}.
\end{equation}
The number of survivors, rather than positivity alone, permits the
passage from a prime cutoff to a bound in terms of the number of prime
divisors. For sufficiently large $k$, if $n$ has at most $k$ distinct
prime divisors, apply
Theorem~\ref{thm:survivors} at a cutoff
$z=A_0 k\log k/\log\log k$, where $A_0$ is a sufficiently large fixed
constant. A progression upper sieve bounds the total number of these
survivors removed by the remaining divisor primes $q>z$ by
$O(k(1+X/\log X))$, where $X=Y/z$.
The lower bound in \eqref{eq:survivor-scale} dominates this loss for
that choice of $A_0$, while $Y\asymp k^2/(\log\log k)^2$.
Section~\ref{sec:assembly} gives the comparison with constants and
handles the bounded range of $k$.

The quantitative estimate also relates to the interval-sieve quantity
studied by Banks, Ford and Tao \cite[Sections~1.3--1.4]{BanksFordTao2023}.
They minimize the survivor count over all choices of forbidden classes
at primes up to $(y/\log y)^{1/2}$ in an interval of length $y$.
At that cutoff a classical lower sieve gives a bound of order
$y\log\log y/(\log y)^2$.
Theorem~\ref{thm:survivors} retains this order of magnitude at the
larger cutoff $z\asymp\sqrt Y\log Y$.
The classes here are arbitrary; no random choice of the forbidden
classes is assumed in this theorem.

\subsection*{Outline of the argument}

The proof has two parts: a reference calculation for a decreasing-prime
expansion, and a comparison with the actual residue classes.
Section~\ref{sec:tree} partitions discarded integers by the least
prime at which they hit a forbidden class. Iteration gives tuples of
strictly decreasing primes.
At even depths the expansion is a lower bound, and at odd depths it
is an upper bound. An admission rule limits the progressions retained
in this expansion so that each is long enough for an absolute upper
sieve. Replacing its small-prime count by progression coordinate
length times $V_0$ defines the reference tree.

The leading reference term is governed by the classical linear-sieve
functions $f,F$ \cite[Section~21.2.2]{MontgomeryVaughanII}, with the
root tending to the critical parameter $2$. The lower exponent cutoff
leaves a boundary contribution on the smaller scale $B^{-2}$.
Each prime path has weight equal to the reciprocal product of its
primes. Replacing reciprocal-prime sums by harmonic integrals $dx/x$
gives a continuous reference process. To evaluate the boundary
contribution, Section~\ref{sec:functions} uses derivative weights of
$f,F$ to normalize its transitions to a probability kernel.
Its regeneration cycles give limiting occupation measures,
which Section~\ref{sec:prime-paths} transfers to prime tuples.
Section~\ref{sec:reference} then evaluates the boundary contribution
and bounds the losses imposed by the admission rule. A positive
margin remains after these losses and the negative contribution from
the root's small displacement below parameter $2$.

The actual small-prime counts agree with their reference values in
long progressions, by the fundamental lemma of sieve theory.
At a short progression, a large relative discrepancy must already be
visible along one of the last few edges of its path.
Section~\ref{sec:boxes} places these paths in boxes in which selected
prime factors vary independently. An inverse estimate in
Section~\ref{sec:inverse} shows that many discrepant edges force
alignment with one of finitely many rationals $A/D$: a prime $p$
aligns when $Da_p\equiv A\pmod p$.

The inverse step is related to the principle that modular
concentration forces algebraic structure, developed by Helfgott and
Venkatesh and by Walsh
\cite{HelfgottVenkatesh2009,Walsh2012}.
Our argument uses partial incidence constraints in a single prime bin
and obtains a rational list fixed for the whole box.
It combines polynomial interpolation with Jarn\'{i}k's
convex-arc lattice-point argument \cite{Jarnik1926} and the
pair-counting principle of Gallagher's larger sieve
\cite{Gallagher1971}.
Variance estimates in Section~\ref{sec:variance} then show that a
rational can account for many discrepant endpoints only when its
effective modulus and intercept are small. The pair-moment calculation
is a truncated form of mean singular-series averaging
\cite{Gallagher1976}; the paper proves the required uniformity and
stability under the changes of measure used here.

For the endpoints remaining after this argument, many earlier prime
factors align with the same rational. Section~\ref{sec:stops} stops
their branches at suitable even nodes. The stop rule depends on prime
bins in a way that permits one aligned prime to vary over a large
subset of its bin. Averaging over that prime gives more exact
survivors than the reference subtree requires. A marked-renewal
argument shows that all but an arbitrarily small amount of the
exceptional path mass reaches one of these stops.

The comparison identity \eqref{eq:stopped-comparison-identity} joins
these parts. Proposition~\ref{prop:reference-margin} supplies a positive
reference term; the inverse and variance estimates leave only the
endpoints of Corollary~\ref{cor:hard-endpoints}; and
Propositions~\ref{prop:stopped-comparison} and~\ref{prop:hard-stopping}
supply the stopped comparison and its coverage.
Section~\ref{sec:assembly} bounds the total remaining discrepancy,
proves Theorem~\ref{thm:survivors}, and deduces Theorem~\ref{thm:main}.

\subsection*{Conventions}

The constants in the proof are absolute after the explicitly stated
parameters have been chosen. They need not be effective. Whenever a
bin width $\xi$ occurs, it is fixed before $z$ tends to infinity.
Constants independent of sufficiently small fixed $\xi$ do not assert
uniformity for a width $\xi(z)$ tending to zero. This order matters
in the rational-list and stopping arguments.

We use $\varphi$ for Euler's totient, $\tau$ for the divisor function,
$\rad(m)$ for the product of the distinct prime divisors of $m$, and
$\bar u$ for a multiplicative inverse in a displayed modulus. All
rationals $A/D$ used for alignment are reduced, with $D>0$.
